\def\xhref#1#2{\href{#1}{\hskip16pt}\llap{#2}}
\def\xxhref#1#2{\href{#1}{\hskip14pt}\llap{#2}}
\makeatletter
\def\@oddhead{\hfill}
\def\@evenhead{\hfill}
\makeatother
   
   
\normalsize
\fontsize{9.6}{11.6}\selectfont
%\parskip4pt plus2pt minus1pt

%\advance\baselineskip by-.3pt

\pznspis{Malin~1 -- co nowego u~spiralnej galaktyki gigantki?}


\pzntyt{Malin~1 -- co nowego u~spiralnej galaktyki gigantki?}

Największa znana do tej pory galaktyka spiralna została odkryta przez astronoma Davida Malina w~1986 roku. 
Na cześć odkrywcy nazwano ją Malin~1. Jest to galaktyka o~niskiej jasności powierzchniowej. Jej jasność – oglądana z~Ziemi – jest mniejsza od jasności otaczającego ją nieba (tzn. w~świetle widzialnym jest prawie przezroczysta). Ramiona spiralne Malin~1 są bardzo rozproszone i~z~jakiegoś powodu galaktyka ta produkuje bardzo mało gwiazd. \marg{Malin 1 znajduje się w~odległości 366 Mpc (czyli 
1,2 miliarda lat świetlnych) od Ziemi. Dla porównania, odległość do galaktyki Andromedy wynosi tylko 2,5 miliona lat świetlnych. Średnica Malin~1 to około 200 kpc, czyli  prawie 660~000~lat świetlnych. Dla porównania, szacowana średnica Drogi Mlecznej to tylko 100 000 lat świetlnych. Średnica naszej Galaktyki jest więc prawie siedem razy mniejsza niż średnica Malin~1.
}

Przez prawie 30 lat od odkrycia Malin~1 naukowcy starali się rozwikłać szereg zagadek związanych z~tym monstrualnym, a~równocześnie prawie przezroczystym, obiektem. \marg{Jak pewnie stali Czytelnicy {\em Delty} pamiętają, metaliczność w~astronomii to pojęcie różne od metaliczności, do jakiej jesteśmy przyzwyczajeni. W~astronomii pojęcie metali odnosi się do wszystkich pierwiastków w~układzie okresowym, które nie są wodorem lub helem.\bigskip} 
Próbowano odpowiedzieć na pytania o~to, jakie były początkowe warunki jej powstania oraz jaka była historia galaktyki. Jak to się stało, że jest ona tak ogromna? Czy tworzą się w~niej nowe gwiazdy? Czy zawiera pył międzygalaktyczny? Jaka jest jej metaliczność?
Dopiero najnowsze obserwacje, uzyskane dzięki wysokiej czułości teleskopu ALMA (o~którym pisaliśmy już wielokrotnie na łamach Delty, np. w~\href{https://www.deltami.edu.pl/2021/07/astronomia-milimetrowa-obserwatorium-alma/}{$\Delta^7\sb{21}$}), pozwoliły na odkrycie części tajemnic skrywanych przez Malin~1.   

Malin~1 jest tak wielką galaktyką, że obserwacje instrumentem MUSE na teleskopie VLT (Very Large Telescope), które odbyły się w~2021~roku i~trwały prawie 1,5~godziny, pozwoliły na zbadanie jedynie 1/4 jej obszaru, w~tym centralnej części galaktyki. Na ich podstawie dr~Junais, pracujący wtedy w~Zakładzie Astrofizyki w~Narodowym Centrum Badań Jądrowych, opublikował badania sugerujące, że Malin~1 posiada silnie rozszerzony dysk gwiazdotwórczy, w~którym promieniowanie pochodzące z~młodych obszarów gwiazdotwórczych jest tłumione przez pył międzygwiazdowy. Było to przełomowe odkrycie, gdyż do tej pory analizowano dane pochodzące z~galaktyk o~niskiej jasności powierzchniowej, przy założeniu, że ilość znajdującego się w~nich pyłu jest zbyt mała, by mieć znaczenie. 
\marg{\includegraphics[width=0.9\linewidth]{Malin1image.png}\\[4pt]
    Rys. 1. Pierwsze zdjęcie Malin~1 z~opisującego jej przypadkowe odkrycie artykułu Bothun, Impey, Malin i~Mould, The Astrophysical Journal, 1987 }
Także mierzona metaliczność Malin~1 okazała się zaskakująca. W~centrum galaktyki jest wyższa od metaliczności Słońca (!), a~następnie opada i~utrzymuje w~miarę stałą wartość w~zewnętrznym dysku.
 
Kilka miesięcy później, wykorzystując w~obserwacjach teleskop ALMA, ten sam zespół naukowy, tym razem pod przewodnictwem dr. Gaspara Galaza z~Instituto de Astrofísica, Pontificia Universidad Católica  w~Santiago w~Chile, zaobserwował emisję cząsteczek tlenku węgla w~centrum Malin~1 (patrz rys. 2), a~także potencjalne emisje w~trzech innych regionach dysku. 
Odkrycie to pozwoliło na oszacowanie masy wodoru molekularnego (który jest gazem umożliwiającym tworzenie się nowych gwiazd w~galaktykach) oraz średnią gęstość powierzchniową gazu molekularnego.

    
Wyniki otrzymane przez zespół (oba artykuły zostały opracowane w~tym samym zespole naukowym) pokazują nie tylko, że Malin~1 bardzo powoli tworzy nowe gwiazdy i~ma ekstremalnie niską ogólną gęstość powierzchniową gazu molekularnego, ale także, że gwiazdy tworzą się tylko w~niektórych miejscach dysku galaktyki. Wyniki analizy sugerują, że galaktyki tak duże i~rozproszone jak Malin~1 tworzą gwiazdy w~wielu miejscach wzdłuż swoich ogromnych ramion spiralnych.
\medskip

 {\includegraphics[trim=0 11 0 0,clip,width=.9\linewidth]{image.png}}\marg[-100]{
      Rys. 2. Detekcja emisji ciągłej i~linii $^{12}CO (1-0)$ za pomocą teleskopu ALMA. Lewy panel przedstawia kolorowy obraz optyczny z~pracy Boissier i~inni (2016). Panele po prawej stronie pokazują powiększone obrazy centrum Malin~1 z~nałożoną detekcją linii $^{12}CO (1-0)$ (panel górny) oraz kontinuum (panel dolny). Czerwona elipsa w~lewym dolnym rogu każdego panelu reprezentuje rozmiar wiązki ALMA. Z~artykułu Galaz i~inni (2024)}

\marg[-2pt]{Tekst opracowany na podstawie artykułów: Junais i~inni ,,MUSE observations of the giant low surface brightness galaxy Malin 1: Numerous HII regions, star formation rate, metallicity, and dust attenuation'' -- opublikowany w~czasopiśmie \textit{Astronomy \& Astrophysics} w~styczniu 2024 r. oraz  Gaspar Galaz i~inni ,,First Detection of Molecular Gas in the Giant Low Surface Brightness Galaxy Malin 1'' -- opublikowany w~czasopiśmie \textit{The Astrophysical Journal Letters} w~listopadzie 2024 r.}
Cierpliwość w~przypadku Malin~1 była niezbędna, gdyż dopiero po 30 latach od jej odkrycia dysponujemy teleskopami pozwalającymi nie tylko na obrazowanie tak słabo świecącej i~ogromnej galaktyki, ale także na analizę jej własności fizycznych przy pomocy obserwacji poszczególnych linii emisyjnych. 
Przygoda z~Malin~1 pewnie dopiero się zaczyna, czekamy na dalsze obserwacje i~analizy, gdyż nadal nie wiemy, w~jaki sposób została uformowana tak wielka galaktyka. 




\bigskip






\rightline{\large%\textit{Anna DURKALEC}%
%\textit{Michał BEJGER}%
\textit{Katarzyna MAŁEK}%
}%
%}

\medskip

%{\scriptsize \rightline
%  {Centrum Astronomiczne im. Mikołaja Kopernika PAN,}
%   \rightline  {Istituto Nazionale di Fisica Nucleare (INFN), Sezione di Ferrara, Włochy}\par} 

%{\scriptsize \rightline{Zakład Astrofizyki, Departament Badań Podstawowych, Narodowe Centrum Badań Jądrowych}}  %A.Durkalec

 \rightline{ \scriptsize  Departament Badań Podstawowych (BP4), Zakład Astrofizyki, Narodowe Centrum Badań Jądrowych}%K.Małek
%\vfill
% \end{dwieszpalty}

\bigskip\medskip
%\normalsize
\input 10-niebo




\endinput